% =========================================================
% frontmatter/portfolio_context.tex
% =========================================================

\chapter*{Portfolio Context and Evidence}
\thispagestyle{fancy}
\addcontentsline{toc}{chapter}{Portfolio Context and Evidence}

\noindent
\textbf{Project Context.}
This work originated from a graduate-level antenna engineering simulation project.
The present document is an independently prepared public portfolio edition intended
to demonstrate analytical antenna design, full-wave electromagnetic simulation,
length tuning, radiation-pattern verification, engineering interpretation, and
technical documentation.

\vspace{0.5em}

\noindent
\textbf{Public Portfolio Edition.}
The original academic report has been reorganized and edited for professional and
educational use. Course-specific grading material, student-identification information,
institutional branding, and administrative content have been removed. This document
does not represent an official publication or endorsement by any university,
instructor, software vendor, employer, or other institution.

\vspace{0.5em}

\noindent
\textbf{Evidence Classification.}
The report combines analytical calculations with CST Studio Suite full-wave simulation.
Analytical evidence includes wavelength and half-wave estimates, the canonical thin
half-wave dipole pattern, and an independent Python calculation of directivity and
half-power beamwidth. Simulated evidence includes CST geometry screenshots,
reflection-coefficient plots, far-field plots, and the exact marker for the final tuned
resonance. No antenna prototype was fabricated, and no calibrated impedance or
antenna-range measurements were conducted. Accordingly, the reported results are
analytical or simulated evidence, not measured performance.

\vspace{0.5em}

\noindent
\textbf{Simulation Scope.}
The preserved source archive contains exported CST figures rather than the native CST
model or raw numerical curve exports. The final tuned resonance is supported by an
exact displayed marker; the preliminary 150~mm and 125~mm resonance values are
estimated from plotted curves. Solver version, mesh statistics, stopping criteria, and
formal mesh- or boundary-convergence records were not preserved. These limitations
are stated explicitly in the numerical-reliability discussion.

\vfill

\begin{center}
    \textit{Prepared as part of the Hossein Molhem Engineering Portfolio.}
\end{center}
